\scanpage{99}

Notons qu'on a une immersion canonique

\[
X\lhook\joinrel\longrightarrow X^{S[i]}
\]

qui ici s'explicite à partir de

\[
\begin{aligned}\tilde X&\longrightarrow\tilde X^{S[i]}\\ x&\longmapsto x+\underline i\,0\end{aligned}
\]

\deleted{\uline{Lemme} Soit $\unsure{\zeta|z|\xi}=x+iy\in\Gamma\tilde X^{S[i]}$, \inserted{i.e. $q(x)=q(y)$, $\varphi(x,y)=0$ $x+\underline i y$ partout non nul} Pour que $\zeta\in\Gamma\tilde U$ i.e. $q(x)=q(y)\in\Gamma G_m$, il f. et s. que $\exists\,\alpha=a+ib\in\deleted{\unsure{G}}\Gamma G_m^{S[\underline i]}$ \inserted{$(a^2+b^2\in\Gamma G_m)$} et $x_0\in\Gamma\tilde X$ ($q(x_0)=0$ $x_0$ \unsure{non nul en} tte fibre \unsure{\ldots}) tels que $\zeta=\alpha x_0$ i.e.}

\[
\deleted{x=ax_0,\ y=bx_0}
\]

En \deleted{fait} on trouve que

\uline{Proposition} L'application inverse de $f:\inserted{U\xrightarrow{\sim}\Sigma}$ (\unsure{déduit|déduite} par passage au quotient de $\tilde f:\tilde U\to\Sigma$, $\tilde f(x+iy)=[x,y]/q(x)$) induit un isomorphisme

\[
\boxed{\Sigma\xrightarrow{\sim}\tilde X^{S[i]}-\tilde X=U}
\]

Ceci équivaut au \deleted{lemme}

\uline{Lemme} Si $S=\operatorname{Spec} k$, $k$ un corps, \inserted{alg clos de car $\ne 2$} alors pour \deleted{\unsure{que}} $x+\underline i y\in\Gamma\tilde X^{S[i]}$ [$x,y$ deux él. de $E$ \inserted{non nuls tous deux} tels que \deleted{$x,y$} \deleted{\unsure{non coll. ou|non colin., soit}, $a+ib$ \unsure{ou non} tels que} \inserted{l'un \unsure{au|ou} \unsure{moins}} i.e. $y=\lambda x$, \deleted{\unsure{ou ait} \unsure{$x\ne0$}} \inserted{avec} $\lambda^2\ne-1$, \inserted{et $q(x)=q(y)$, $\varphi(x,y)=0$}] \deleted{\unsure{soit|alors}} \unsure{;|,} on a $x+iy\notin\tilde U^{\deleted{\unsure{S[i]}}}$ i.e. $q(x)=q(y)=0$ sss $x+\underline i y=(a+\underline i b)x_0$, avec $x_0\in\Gamma\tilde X$ i.e. $\unsure{\cdot}x_0\in E^{*}$, $q(x_0)=0$, \deleted{$a+b$} $a+\underline i b\in\Gamma G_m^{S[i]}$ i.e $a^2+b^2\ne0$.

% archive sheet number: 98
